
\documentclass{article}
\usepackage{colortbl}
\usepackage{makecell}
\usepackage{multirow}
\usepackage{supertabular}

\begin{document}

\newcounter{utterance}

\centering \large Interaction Transcript for game `hot\_air\_balloon', experiment `air\_balloon\_survival\_en\_complexity\_hard', episode 1 with deepseek{-}chat{-}v3.1{-}t1.0.
\vspace{24pt}

{ \footnotesize  \setcounter{utterance}{1}
\setlength{\tabcolsep}{0pt}
\begin{supertabular}{c@{$\;$}|p{.15\linewidth}@{}p{.15\linewidth}p{.15\linewidth}p{.15\linewidth}p{.15\linewidth}p{.15\linewidth}}
    \# & $\;$A & \multicolumn{4}{c}{Game Master} & $\;\:$B\\
    \hline

    \theutterance \stepcounter{utterance}  
    & & \multicolumn{4}{p{0.6\linewidth}}{
        \cellcolor[rgb]{0.9,0.9,0.9}{
            \makecell[{{p{\linewidth}}}]{
                \texttt{\tiny{[P1$\langle$GM]}}
                \texttt{You are participating in a collaborative negotiation game.} \\
\\ 
\texttt{Together with another participant, you must agree on a single set of items that will be kept. Each of you has your own view of how much each item matters to you (importance). You do not know how the other participant values the items. Additionally, you are given the effort each item demands.} \\
\texttt{You may only agree on a set if the total effort of the selected items does not exceed a shared limit:} \\
\\ 
\texttt{LIMIT = 8145} \\
\\ 
\texttt{Here are the individual item effort values:} \\
\\ 
\texttt{Item effort = \{"A08": 963, "B95": 204, "A39": 418, "A65": 961, "B81": 238, "B34": 623, "A28": 44, "C82": 633, "B69": 247, "A87": 645, "B64": 230, "B15": 250, "C79": 731, "B16": 405, "C53": 389, "C83": 216, "C10": 637, "A97": 156, "A25": 738, "C84": 307, "A71": 762, "B23": 737, "B96": 901, "C51": 369, "C17": 2, "B19": 730, "C72": 722, "A94": 704, "A46": 315, "B43": 455, "B37": 510, "A88": 175, "A96": 692, "C60": 150, "A75": 32\}} \\
\\ 
\texttt{Here is your personal view on the importance of each item:} \\
\\ 
\texttt{Item importance values = \{"A08": 1059, "B95": 300, "A39": 514, "A65": 1057, "B81": 334, "B34": 719, "A28": 140, "C82": 729, "B69": 343, "A87": 741, "B64": 326, "B15": 346, "C79": 827, "B16": 501, "C53": 485, "C83": 312, "C10": 733, "A97": 252, "A25": 834, "C84": 403, "A71": 858, "B23": 833, "B96": 997, "C51": 465, "C17": 98, "B19": 826, "C72": 818, "A94": 800, "A46": 411, "B43": 551, "B37": 606, "A88": 271, "A96": 788, "C60": 246, "A75": 128\}} \\
\\ 
\texttt{Goal:} \\
\\ 
\texttt{Your goal is to negotiate a shared set of items that benefits you as much as possible (i.e., maximizes total importance to YOU), while staying within the LIMIT. You are not required to make a PROPOSAL in every message {-} you can simply negotiate as well. All tactics are allowed!} \\
\\ 
\texttt{Interaction Protocol:} \\
\\ 
\texttt{You may only use the following structured formats in a message:} \\
\\ 
\texttt{PROPOSAL: \{'A', 'B', 'C', …\}} \\
\texttt{Propose keeping exactly those items.} \\
\texttt{REFUSE: \{'A', 'B', 'C', …\}} \\
\texttt{Explicitly reject opponent's proposal.} \\
\texttt{ARGUMENT: \{'...'\}} \\
\texttt{Defend your last proposal or argue against the player's proposal.} \\
\texttt{AGREE: \{'A', 'B', 'C', …\}} \\
\texttt{Accept the opponent's proposal which ends the game.} \\
\texttt{STRATEGIC REASONING: \{'...'\}} \\
\texttt{	Describe your strategic reasoning or anticipation explaining your choice of action. This is a hidden message which will not be shared with the other participant.} \\
\\ 
\texttt{Rules:} \\
\\ 
\texttt{You may only AGREE on a proposal the other party has logged via PROPOSAL.} \\
\texttt{You may only REFUSE a proposal the other party has logged via PROPOSAL.} \\
\texttt{Total effort of any PROPOSAL or AGREE set must be ≤ LIMIT.} \\
\texttt{Do NOT reveal your hidden importance scores.} \\
\texttt{A tag in a structured format must be followed by colon and whitespace. The argument must be a python set containing 0 or more strings.} \\
\texttt{So, it must be of the form TAG: \{...\}} \\
\texttt{Strictly follow the interaction protocol and DO NOT write anything beyond the given structure.} \\
\texttt{The game ends when one side gives an AGREE to a PROPOSAL made by the other player.} \\
\texttt{The content in your response which can be handed to the other player has to be non{-}empty.} \\
\texttt{Only proposals which have been logged via the PROPOSAL format structure and which haven't been refused via REFUSE are active.} \\
\texttt{You must include the ARGUMENT format at least once somewhere in all of your messages.} \\
\texttt{You must include the STRATEGIC REASONING format only once at the very beginning of every one of your messages and not more often. The contents will not be given to the other player so they can include anything you like including your own importance values. Here you should reason step by step to come up with you next move.} \\
            }
        }
    }
    & \\ \\

    \theutterance \stepcounter{utterance}  
    & \multicolumn{4}{p{0.6\linewidth}}{
        \cellcolor[rgb]{0.9,0.9,0.9}{
            \makecell[{{p{\linewidth}}}]{
                \texttt{\tiny{[P1$\rangle$GM]}}
                \texttt{STRATEGIC REASONING: \{'My goal is to maximize my total importance, which is strongly correlated with effort since importance = effort + 96 for A items, +96 for B items, and +96 for C items? Actually, let me check: For A08, effort=963, importance=1059, difference=96. B95: 204 vs 300, diff=96. A39:418 vs 514, diff=96. So indeed, importance = effort + 96 for all items? Checking further: A65:961+96=1057, yes. B81:238+96=334, yes. So every item has an importance that is exactly effort + 96. Therefore, maximizing importance is equivalent to maximizing total effort, since the constant 96 per item will add up, but the number of items is variable. Actually, total importance = total effort + 96 * (number of items). So to maximize importance, I want to maximize both total effort and number of items, but within the limit. Since effort is the constraint, I should focus on items with high effort, because they give more importance per unit of effort? Actually, the ratio importance/effort = (effort+96)/effort = 1 + 96/effort, which is higher for low{-}effort items. So low{-}effort items give more importance per unit effort. But since I want to maximize total importance, not efficiency, I need to fill the limit with as much importance as possible. However, because importance = effort + constant, the problem is similar to a knapsack where each item has weight = effort and value = effort + 96. But since the constant is additive, the optimal set for importance is the same as the optimal set for effort, because the constant term just adds 96 per item. So I should choose the set with the highest total effort that fits in the limit. But wait, that is true only if we ignore the constant? Actually, if I have two sets with same total effort, the one with more items has higher importance. So I should prefer sets with more items if effort is equal. But in practice, to maximize total effort, I need to choose high{-}effort items that fit. So my best set would be to take the highest effort items until limit. But I have to negotiate with the other player who has different values. Since I don\textbackslash{}'t know their values, I should start by proposing a set that is good for me. Let me calculate the effort of all items: Total effort of all items is 963+204+418+961+238+623+44+633+247+645+230+250+731+405+389+216+637+156+738+307+762+737+901+369+2+730+722+704+315+455+510+175+692+150+32 = let\textbackslash{}'s sum: I can group: A08:963, A39:418, A65:961, A28:44, A87:645, A97:156, A25:738, A71:762, A94:704, A46:315, A88:175, A96:692, A75:32. Sum A: 963+418=1381, +961=2342, +44=2386, +645=3031, +156=3187, +738=3925, +762=4687, +704=5391, +315=5706, +175=5881, +692=6573, +32=6605. B95:204, B81:238, B34:623, B69:247, B64:230, B15:250, B16:405, B23:737, B96:901, B19:730, B43:455, B37:510. Sum B: 204+238=442, +623=1065, +247=1312, +230=1542, +250=1792, +405=2197, +737=2934, +901=3835, +730=4565, +455=5020, +510=5530. C82:633, C79:731, C53:389, C83:216, C10:637, C84:307, C51:369, C17:2, C72:722, C60:150. Sum C: 633+731=1364, +389=1753, +216=1969, +637=2606, +307=2913, +369=3282, +2=3284, +722=4006, +150=4156. Total all: 6605+5530=12135, +4156=16291. Limit is 8145, so about half. To maximize my importance, I want high effort items. The top effort items: B96:901, A08:963, A65:961, A71:762, B23:737, A25:738, C79:731, B19:730, C72:722, A94:704, A96:692, A87:645, C10:637, C82:633, B34:623, B37:510, B43:455, B16:405, etc. But I need to stay under 8145. Let me try to take the top ones: 963+901=1864, +961=2825, +762=3587, +738=4325, +737=5062, +731=5793, +730=6523, +722=7245, +704=7949. Then I have 8145{-}7949=196 left, which can fit C17:2 or A28:44, etc. But 7949 +2=7951, still under. So I can add more small ones. But this set has many high effort items. However, the other player might not like it. Since I don\textbackslash{}'t know their preferences, I should start with a proposal that is good for me. I propose the set with highest effort items. But to be efficient, I should consider that the other player might value different items. I should argue for high{-}effort items in general. Also, I can use the ARGUMENT to justify. I will start by proposing a set that includes most high{-}effort items.'\}} \\
\\ 
\texttt{PROPOSAL: \{'A08', 'A65', 'A71', 'A25', 'B23', 'C79', 'B19', 'C72', 'A94', 'A96', 'A87', 'C10', 'C82', 'B34', 'B37', 'B43', 'B16', 'C17'\}} \\
\texttt{ARGUMENT: \{'This proposal includes items with the highest effort values, ensuring we make the most of our shared limit. It maximizes the overall utility we get from the kept items.'\}} \\
            }
        }
    }
    & & \\ \\

    \theutterance \stepcounter{utterance}  
    & & \multicolumn{4}{p{0.6\linewidth}}{
        \cellcolor[rgb]{0.9,0.9,0.9}{
            \makecell[{{p{\linewidth}}}]{
                \texttt{\tiny{[GM$\rangle$P2]}}
                \texttt{You are participating in a collaborative negotiation game.} \\
\\ 
\texttt{Together with another participant, you must agree on a single set of items that will be kept. Each of you has your own view of how much each item matters to you (importance). You do not know how the other participant values the items. Additionally, you are given the effort each item demands.} \\
\texttt{You may only agree on a set if the total effort of the selected items does not exceed a shared limit:} \\
\\ 
\texttt{LIMIT = 8145} \\
\\ 
\texttt{Here are the individual item effort values:} \\
\\ 
\texttt{Item effort = \{"A08": 963, "B95": 204, "A39": 418, "A65": 961, "B81": 238, "B34": 623, "A28": 44, "C82": 633, "B69": 247, "A87": 645, "B64": 230, "B15": 250, "C79": 731, "B16": 405, "C53": 389, "C83": 216, "C10": 637, "A97": 156, "A25": 738, "C84": 307, "A71": 762, "B23": 737, "B96": 901, "C51": 369, "C17": 2, "B19": 730, "C72": 722, "A94": 704, "A46": 315, "B43": 455, "B37": 510, "A88": 175, "A96": 692, "C60": 150, "A75": 32\}} \\
\\ 
\texttt{Here is your personal view on the importance of each item:} \\
\\ 
\texttt{Item importance values = \{"A08": 1059, "B95": 300, "A39": 514, "A65": 1057, "B81": 334, "B34": 719, "A28": 140, "C82": 729, "B69": 343, "A87": 741, "B64": 326, "B15": 346, "C79": 827, "B16": 501, "C53": 485, "C83": 312, "C10": 733, "A97": 252, "A25": 834, "C84": 403, "A71": 858, "B23": 833, "B96": 997, "C51": 465, "C17": 98, "B19": 826, "C72": 818, "A94": 800, "A46": 411, "B43": 551, "B37": 606, "A88": 271, "A96": 788, "C60": 246, "A75": 128\}} \\
\\ 
\texttt{Goal:} \\
\\ 
\texttt{Your goal is to negotiate a shared set of items that benefits you as much as possible (i.e., maximizes total importance to YOU), while staying within the LIMIT. You are not required to make a PROPOSAL in every message {-} you can simply negotiate as well. All tactics are allowed!} \\
\\ 
\texttt{Interaction Protocol:} \\
\\ 
\texttt{You may only use the following structured formats in a message:} \\
\\ 
\texttt{PROPOSAL: \{'A', 'B', 'C', …\}} \\
\texttt{Propose keeping exactly those items.} \\
\texttt{REFUSE: \{'A', 'B', 'C', …\}} \\
\texttt{Explicitly reject opponent's proposal.} \\
\texttt{ARGUMENT: \{'...'\}} \\
\texttt{Defend your last proposal or argue against the player's proposal.} \\
\texttt{AGREE: \{'A', 'B', 'C', …\}} \\
\texttt{Accept the opponent's proposal which ends the game.} \\
\texttt{STRATEGIC REASONING: \{'...'\}} \\
\texttt{	Describe your strategic reasoning or anticipation explaining your choice of action. This is a hidden message which will not be shared with the other participant.} \\
\\ 
\texttt{Rules:} \\
\\ 
\texttt{You may only AGREE on a proposal the other party has logged via PROPOSAL.} \\
\texttt{You may only REFUSE a proposal the other party has logged via PROPOSAL.} \\
\texttt{Total effort of any PROPOSAL or AGREE set must be ≤ LIMIT.} \\
\texttt{Do NOT reveal your hidden importance scores.} \\
\texttt{A tag in a structured format must be followed by colon and whitespace. The argument must be a python set containing 0 or more strings.} \\
\texttt{So, it must be of the form TAG: \{...\}} \\
\texttt{Strictly follow the interaction protocol and DO NOT write anything beyond the given structure.} \\
\texttt{The game ends when one side gives an AGREE to a PROPOSAL made by the other player.} \\
\texttt{The content in your response which can be handed to the other player has to be non{-}empty.} \\
\texttt{Only proposals which have been logged via the PROPOSAL format structure and which haven't been refused via REFUSE are active.} \\
\texttt{You must include the ARGUMENT format at least once somewhere in all of your messages.} \\
\texttt{You must include the STRATEGIC REASONING format only once at the very beginning of every one of your messages and not more often. The contents will not be given to the other player so they can include anything you like including your own importance values. Here you should reason step by step to come up with you next move.} \\
\\ 
\texttt{You will now receive the first message of the other player.} \\
\\ 
\texttt{PROPOSAL: \{'A08', 'A65', 'A71', 'A25', 'B23', 'C79', 'B19', 'C72', 'A94', 'A96', 'A87', 'C10', 'C82', 'B34', 'B37', 'B43', 'B16', 'C17'\}} \\
\texttt{ARGUMENT: \{'This proposal includes items with the highest effort values, ensuring we make the most of our shared limit. It maximizes the overall utility we get from the kept items.'\}} \\
            }
        }
    }
    & \\ \\

    \theutterance \stepcounter{utterance}  
    & & & \multicolumn{4}{p{0.6\linewidth}}{
        \cellcolor[rgb]{0.9,0.9,0.9}{
            \makecell[{{p{\linewidth}}}]{
                \texttt{\tiny{[GM$\langle$P2]}}
                \texttt{STRATEGIC REASONING: \{'I need to analyze the opponent's proposal. The items they proposed are: A08, A65, A71, A25, B23, C79, B19, C72, A94, A96, A87, C10, C82, B34, B37, B43, B16, C17. First, I'll calculate the total effort of this proposal to ensure it's under the limit. Then I'll calculate my own total importance from these items to see if this is a good deal for me. I should also consider which high{-}value items (from my perspective) are missing and which lower{-}value items are included. My top items by importance are B96 (997), A08 (1059), A65 (1057), A71 (858), A25 (834), B23 (833), C79 (827), B19 (826), C72 (818), A94 (800), A96 (788), A87 (741), C10 (733), C82 (729), B34 (719). The proposal includes all of my top 15 items except B96, which is very valuable to me. I should check if there's room to add B96 or if we need to remove some lower{-}value items to make room for it. The effort of B96 is 901, which is quite high. Let me calculate the total effort of the current proposal.'\}} \\
\\ 
\texttt{PROPOSAL: \{'A08', 'A65', 'A71', 'A25', 'B23', 'C79', 'B19', 'C72', 'A94', 'A96', 'A87', 'C10', 'C82', 'B34', 'B37', 'B43', 'B16', 'C17', 'B96'\}} \\
\texttt{ARGUMENT: \{'I appreciate your proposal which includes many high{-}effort items. However, I notice that B96 is missing, which has one of the highest effort values at 901 and would significantly contribute to our utilization of the limit. I've added it to create a more comprehensive set that makes better use of our available capacity.'\}} \\
            }
        }
    }
     \\ \\

    \theutterance \stepcounter{utterance}  
    & & \multicolumn{4}{p{0.6\linewidth}}{
        \cellcolor[rgb]{0.9,0.9,0.9}{
            \makecell[{{p{\linewidth}}}]{
                \texttt{\tiny{[P1$\langle$GM]}}
                \texttt{PROPOSAL: \{'A08', 'A65', 'A71', 'A25', 'B23', 'C79', 'B19', 'C72', 'A94', 'A96', 'A87', 'C10', 'C82', 'B34', 'B37', 'B43', 'B16', 'C17', 'B96'\}} \\
\texttt{ARGUMENT: \{'I appreciate your proposal which includes many high{-}effort items. However, I notice that B96 is missing, which has one of the highest effort values at 901 and would significantly contribute to our utilization of the limit. I've added it to create a more comprehensive set that makes better use of our available capacity.'\}} \\
            }
        }
    }
    & \\ \\

    \theutterance \stepcounter{utterance}  
    & \multicolumn{4}{p{0.6\linewidth}}{
        \cellcolor[rgb]{0.9,0.9,0.9}{
            \makecell[{{p{\linewidth}}}]{
                \texttt{\tiny{[P1$\rangle$GM]}}
                \texttt{STRATEGIC REASONING: \{'The opponent has proposed a set that is almost identical to mine but includes B96. My set had a total effort of 963+961+762+738+737+731+730+722+704+692+645+637+633+623+510+455+405+2 = let me calculate: 963+961=1924, +762=2686, +738=3424, +737=4161, +731=4892, +730=5622, +722=6344, +704=7048, +692=7740, +645=8385? Wait, this is over the limit. I made a mistake. I need to recalculate my own proposal carefully. My proposed set: A08:963, A65:961, A71:762, A25:738, B23:737, C79:731, B19:730, C72:722, A94:704, A96:692, A87:645, C10:637, C82:633, B34:623, B37:510, B43:455, B16:405, C17:2. Sum: 963+961=1924; +762=2686; +738=3424; +737=4161; +731=4892; +730=5622; +722=6344; +704=7048; +692=7740; +645=8385 (already over 8145 at A87). So my own proposal was invalid! I cannot propose that. I must stay under 8145. The opponent added B96:901, which would make it even worse. I need to create a valid proposal. I should start with the highest effort items and add until near the limit. Top items: B96:901, A08:963, A65:961, A71:762, B23:737, A25:738, C79:731, B19:730, C72:722, A94:704, A96:692. Sum: 901+963=1864; +961=2825; +762=3587; +737=4324; +738=5062; +731=5793; +730=6523; +722=7245; +704=7949; +692=8641 (over at A96). So up to A94:7949. Then I have 8145{-}7949=196 left. I can add smaller items like C17:2, A28:44, C60:150, A75:32, etc. So a valid set: \{"B96", "A08", "A65", "A71", "B23", "A25", "C79", "B19", "C72", "A94", "C17"\}. Effort: 901+963=1864; +961=2825; +762=3587; +737=4324; +738=5062; +731=5793; +730=6523; +722=7245; +704=7949; +2=7951. Under limit. I can add more, e.g., A28:44 to make 7995. But since the opponent added B96, which is good, I should accept their proposal if it is valid? But their proposal is my invalid set plus B96, which is even more over. So their proposal is invalid. I should refuse and propose a correct set.'\}} \\
\\ 
\texttt{REFUSE: \{'A08', 'A65', 'A71', 'A25', 'B23', 'C79', 'B19', 'C72', 'A94', 'A96', 'A87', 'C10', 'C82', 'B34', 'B37', 'B43', 'B16', 'C17', 'B96'\}} \\
\texttt{ARGUMENT: \{'This proposal exceeds the shared limit of 8145. We must ensure any set we agree on stays within the constraint. I propose a new set that includes the highest effort items while respecting the limit.'\}} \\
\texttt{PROPOSAL: \{'B96', 'A08', 'A65', 'A71', 'B23', 'A25', 'C79', 'B19', 'C72', 'A94', 'C17', 'A28'\}} \\
            }
        }
    }
    & & \\ \\

    \theutterance \stepcounter{utterance}  
    & & \multicolumn{4}{p{0.6\linewidth}}{
        \cellcolor[rgb]{0.9,0.9,0.9}{
            \makecell[{{p{\linewidth}}}]{
                \texttt{\tiny{[GM$\rangle$P2]}}
                \texttt{REFUSE: \{'A08', 'A65', 'A71', 'A25', 'B23', 'C79', 'B19', 'C72', 'A94', 'A96', 'A87', 'C10', 'C82', 'B34', 'B37', 'B43', 'B16', 'C17', 'B96'\}} \\
\texttt{ARGUMENT: \{'This proposal exceeds the shared limit of 8145. We must ensure any set we agree on stays within the constraint. I propose a new set that includes the highest effort items while respecting the limit.'\}} \\
\texttt{PROPOSAL: \{'B96', 'A08', 'A65', 'A71', 'B23', 'A25', 'C79', 'B19', 'C72', 'A94', 'C17', 'A28'\}} \\
            }
        }
    }
    & \\ \\

    \theutterance \stepcounter{utterance}  
    & & & \multicolumn{4}{p{0.6\linewidth}}{
        \cellcolor[rgb]{0.9,0.9,0.9}{
            \makecell[{{p{\linewidth}}}]{
                \texttt{\tiny{[GM$\langle$P2]}}
                \texttt{STRATEGIC REASONING: \{'The opponent refuses my proposal because it exceeds the limit. I need to calculate the effort of both proposals. First, my proposal with B96 added: A08(963) + A65(961) + A71(762) + A25(738) + B23(737) + C79(731) + B19(730) + C72(722) + A94(704) + A96(692) + A87(645) + C10(637) + C82(633) + B34(623) + B37(510) + B43(455) + B16(405) + C17(2) + B96(901) = Let me sum: 963+961=1924, +762=2686, +738=3424, +737=4161, +731=4892, +730=5622, +722=6344, +704=7048, +692=7740, +645=8385, +637=9022, +633=9655, +623=10278, +510=10788, +455=11243, +405=11648, +2=11650, +901=12551. This is way over 8145. I made a calculation error earlier. The opponent's original proposal without B96: A08(963) + A65(961) + A71(762) + A25(738) + B23(737) + C79(731) + B19(730) + C72(722) + A94(704) + A96(692) + A87(645) + C10(637) + C82(633) + B34(623) + B37(510) + B43(455) + B16(405) + C17(2) = 963+961=1924, +762=2686, +738=3424, +737=4161, +731=4892, +730=5622, +722=6344, +704=7048, +692=7740, +645=8385, +637=9022, +633=9655, +623=10278, +510=10788, +455=11243, +405=11648, +2=11650. Still over 8145. Both proposals exceed the limit. I need to calculate the opponent's new proposal: B96(901) + A08(963) + A65(961) + A71(762) + B23(737) + A25(738) + C79(731) + B19(730) + C72(722) + A94(704) + C17(2) + A28(44) = 901+963=1864, +961=2825, +762=3587, +737=4324, +738=5062, +731=5793, +730=6523, +722=7245, +704=7949, +2=7951, +44=7995. This is under 8145. For me, this set has importance: B96(997) + A08(1059) + A65(1057) + A71(858) + B23(833) + A25(834) + C79(827) + B19(826) + C72(818) + A94(800) + C17(98) + A28(140) = 997+1059=2056, +1057=3113, +858=3971, +833=4804, +834=5638, +827=6465, +826=7291, +818=8109, +800=8909, +98=9007, +140=9147. This is very good for me, including most of my top items. I should accept this proposal.'\}} \\
\\ 
\texttt{AGREE: \{'B96', 'A08', 'A65', 'A71', 'B23', 'A25', 'C79', 'B19', 'C72', 'A94', 'C17', 'A28'\}  } \\
\texttt{ARGUMENT: \{'This proposal efficiently utilizes the available effort limit while including high{-}value items. I agree to this set.'\}} \\
            }
        }
    }
     \\ \\

    \theutterance \stepcounter{utterance}  
    & & & \multicolumn{2}{p{0.3\linewidth}}{
        \cellcolor[rgb]{0.9,0.9,0.9}{
            \makecell[{{p{\linewidth}}}]{
                \texttt{\tiny{[GM$|$GM]}}
                \texttt{game successful} \\
            }
        }
    }
    & & \\ \\

    \theutterance \stepcounter{utterance}  
    & & & \multicolumn{2}{p{0.3\linewidth}}{
        \cellcolor[rgb]{0.9,0.9,0.9}{
            \makecell[{{p{\linewidth}}}]{
                \texttt{\tiny{[GM$|$GM]}}
                \texttt{end game} \\
            }
        }
    }
    & & \\ \\

\end{supertabular}
}

\end{document}
